% -------------
% Preambule pro DRILY (cvičné úlohy ve stylu statnic) - MFF UK, UI / Strojové učení
% Sazba: pdflatex (UTF-8, čeština). Samostatné PDF na úroveň (1/2/3).
% -------------
\usepackage[T1]{fontenc}
\usepackage[utf8]{inputenc}
\usepackage{lmodern}
\usepackage[czech]{babel}
\usepackage[a4paper,margin=2.3cm]{geometry}

\usepackage{amsmath,amssymb,mathtools}
\usepackage{textcomp}
\usepackage[nopatch=all]{microtype}
\usepackage{xcolor}
\usepackage{enumitem}
\usepackage{titlesec}
\usepackage[most]{tcolorbox}
\usepackage{booktabs}
\usepackage{listings}
\usepackage{hyperref}

% - seznamy -
\setlist{leftmargin=1.4em, itemsep=2pt, topsep=3pt, parsep=0pt}

% - barvy -
\definecolor{mffblue}{RGB}{27,54,93}
\definecolor{mffaccent}{RGB}{0,90,156}
\definecolor{ulohagray}{RGB}{70,70,75}
\definecolor{resgreen}{RGB}{30,110,60}
\definecolor{gapred}{RGB}{150,30,30}
\definecolor{calorange}{RGB}{180,110,0}
\definecolor{codebg}{RGB}{245,245,245}

\titleformat{\section}{\Large\bfseries\color{mffblue}}{\thesection}{0.6em}{}
\titleformat{\subsection}{\large\bfseries\color{mffaccent}}{\thesubsection}{0.6em}{}
\titlespacing*{\section}{0pt}{1.4em}{0.6em}
\setcounter{secnumdepth}{2}
\setcounter{tocdepth}{1}

\hypersetup{colorlinks=true, linkcolor=mffaccent, urlcolor=mffaccent,
  pdftitle={Cvičné úlohy ke státnicím - Informatika / UI / Strojové učení},
  pdfauthor={Viktor Číhal}, bookmarksnumbered=true}

\renewcommand{\arraystretch}{1.2}
\setlength{\parindent}{0pt}
\setlength{\parskip}{5pt plus 2pt minus 1pt}
\emergencystretch=3em
\hbadness=10000

% - kód (C#); komentáře držíme bez diakritiky kvuli listings+utf8 -
\lstdefinestyle{cs}{
  basicstyle=\ttfamily\footnotesize, backgroundcolor=\color{codebg},
  keywordstyle=\color{mffaccent}\bfseries, commentstyle=\color{resgreen!70!black}\itshape,
  stringstyle=\color{gapred}, showstringspaces=false, columns=fullflexible,
  breaklines=true, frame=single, framerule=0pt, rulecolor=\color{codebg},
  xleftmargin=6pt, xrightmargin=6pt, aboveskip=4pt, belowskip=4pt,
  literate={ě}{{\v e}}1 {š}{{\v s}}1 {č}{{\v c}}1 {ř}{{\v r}}1 {ž}{{\v z}}1
           {á}{{\'a}}1 {í}{{\'\i}}1 {é}{{\'e}}1 {ý}{{\'y}}1 {ů}{{\r u}}1 {ú}{{\'u}}1,
  language=[Sharp]C}
\lstset{style=cs}

% =====================================================================
% Prostředí pro ÚLOHU a ŘEŠENÍ ve stylu statnic.
% =====================================================================

% čítač úloh
\newcounter{uloha}
\renewcommand{\theuloha}{\arabic{uloha}}

% \uloha{okruh}{téma} -> nadpis úlohy ve stylu zadání statnic
% okruh: "společná matematika" / "společná informatika" / "specializace UI-SU"
\newcommand{\ulohahead}[2]{%
  \refstepcounter{uloha}%
  \par\medskip\noindent
  {\large\bfseries\color{ulohagray}Úloha \theuloha.}\quad
  {\small\colorbox{ulohagray!12}{\textbf{#1}}}\quad{\bfseries #2}\par
  \vspace{2pt}}

% zadání: lehce rámovaný blok jako na zkoušce
\newtcolorbox{zadani}{enhanced, breakable, sharp corners=south,
  colback=ulohagray!5, colframe=ulohagray, boxrule=0pt, leftrule=3pt,
  left=8pt, right=8pt, top=4pt, bottom=5pt, before skip=4pt, after skip=4pt}

% podotázky se značkou bodu (kolik bodu daná část typicky nese)
\newcommand{\bod}[1]{\tikz[baseline=-0.55ex]\node[circle,fill=mffaccent,text=white,
  inner sep=1.3pt,font=\bfseries\scriptsize]{#1};}
% bez tikz (rychlejší): kroužek přes textový symbol
\renewcommand{\bod}[1]{{\footnotesize\colorbox{mffaccent}{\textcolor{white}{\bfseries #1\,b}}}}

% řešení
\newtcolorbox{reseni}{enhanced, breakable, colback=resgreen!6, colframe=resgreen,
  boxrule=0.6pt, arc=1.5pt, left=7pt, right=7pt, top=3pt, bottom=4pt,
  before skip=4pt, after skip=8pt, fonttitle=\bfseries\small, coltitle=resgreen,
  title={Řešení}}

% kalibrační poznámka: jak tato úloha sytí podlahu (common>=5 / spec>=7)
\newcommand{\kalibrace}[1]{\par\smallskip\noindent
  {\footnotesize\textcolor{calorange}{\textbf{Kalibrace:}}~\textit{#1}}\par\smallskip}

% past / častá chyba
\newcommand{\past}[1]{\par\smallskip\noindent
  \textbf{\color{gapred}Pozor (častá chyba):}~#1\par}

% info box
\newtcolorbox{infobox}[1][]{enhanced, breakable, colback=mffaccent!5, colframe=mffaccent,
  boxrule=0.7pt, left=6pt, right=6pt, top=4pt, bottom=4pt, arc=2pt,
  before skip=6pt, after skip=6pt, fonttitle=\bfseries, #1}
